%!TEX root=../main.tex

\section{Relaxing Equalized Odds to Preserve Calibration}
\label{sec:relaxation}

In this section, we show that a substantially simplified notion of Equalized Odds
is compatible with calibration.
We introduce a general relaxation that seeks
to satisfy a \emph{single equal-cost constraint} while maintaining calibration for each group $G_t$.
%The equal-cost constraints we are concerned
%with are based on a class of cost functions defined below.
%
We begin with the observation that Equalized Odds
sets constraints to equalize
false-positives $\fp{h_t}$ and false-negatives $\fn{h_t}$. To capture and
generalize this, we define a
\emph{cost function} \fairconstname{} to be a linear function in \fp{h_t} and
\fn{h_t} with arbitrary dependence on the group's base rate $\br_t$.
More formally, a cost function for group $G_t$ is
\begin{equation}
  \fairconst{h_t}{t} = a_t \fp{h_t} + b_t \fn{h_t}
  \label{eqn:cost_general}
\end{equation}
where $a_{t}$ and $b_{t}$ are non-negative constants that are specific to each group (and
thus may depend on $\br_t$): see \autoref{fig:general-disparity}. We also make
the assumption that for any $\br_t$,
at least one of $a_t$ and $b_t$ is nonzero, meaning
$\fairconst{h_t}{t} = 0$ if and only if $\fp{h_t} = \fn{h_t} = 0$.\footnote{
  By calibration, we cannot have one of $\fp{h_t} = 0$ or $\fn{h_t} = 0$
  without the other, see \autoref{fig:calibration}.
}
This class of cost functions encompasses a
variety of scenarios.
As an example, imagine an application in which the equal false-positive
condition is essential but not the
false-negative condition.
Such a scenario may arise in our recidivism-prediction example, if we require
that non-repeat offenders of any race are not disproportionately labeled as high risk.
%is more likely to experience a missed detection of disease.
If we plot the set of calibrated classifiers $\mathcal{H}^*_1$ and $\mathcal{H}^*_2$
on the false-positive/false-negative plane, we can see that ensuring the
false-positive condition requires finding classifiers $h_1 \in \mathcal{H}^*_1$
and $h_2 \in \mathcal{H}^*_2$ that fall on the same vertical line
(\autoref{fig:two-of-three-fp}). Conversely, if we instead choose to satisfy only
the false-negative condition, we would find classifiers $h_1$ and $h_2$ that
fall on the same horizontal (\autoref{fig:two-of-three-fn}).
%
Finally, if both false-positive and false-negative errors
incur a negative cost on the individual, we may choose to equalize a weighted
combination of the error rates \cite{berk2016primer,berk2017fairness,chouldechova2017fair},
which can be graphically described by the classifiers lying on a convex and negatively-sloped level set
(\autoref{fig:general-disparity}).
%We can express this cost as a weighted combination
%of the generalized error rates
%with a function of the form $\fairconst{h_t}{t} = a_t \fp{h_t} + b_t \fn{h_t}$,
%We can  $a_t$ and $b_t$ are constants that weight each error type which may depend on $\mu_t$
With these definitions,
we can formally define our relaxation:
%
\begin{definition}[Relaxed Equalized Odds with Calibration]
  \label{def:fair}
  Given a cost function $\fairconstname$ of the form in $\eqref{eqn:cost_general}$,
  classifiers $h_1$ and $h_2$ achieve \emph{Relaxed Equalized Odds with
  Calibration} for groups $G_1$ and $G_2$ if both classifiers
  are calibrated and satisfy the constraint $\fairconst{h_1}{1} =
  \fairconst{h_2}{2}$.
\end{definition}
%
It is worth noting that, for calibrated classifiers, an increase in cost
strictly corresponds to an increase in both the false-negative and
false-positive rate. This can be interpreted graphically, as the level-order cost
curves lie further away from the origin as cost increases (\autoref{fig:levelorder}). In other words,
the cost function can always be used as a proxy for either error rate.\footnote{
  This holds even for approximately calibrated classifiers ---  see \autoref{sec:cost_and_error}.
}

\paragraph{Feasibility.}
It is easy to see that \autoref{def:fair} is always satisfiable -- in
Figures~\ref{fig:two-of-three-fp}, \ref{fig:two-of-three-fn}, and \ref{fig:general-disparity}
we see that there are many such solutions that would lie on a given level-order
cost curve while maintaining calibration, including the case in which both
classifiers are perfect.  In
practice, however, not all classifiers are achievable. For the rest of
the paper, we will assume that we have access to ``optimal'' (but possibly
discriminatory) calibrated
classifiers $h_1$ and $h_2$ such that, due to whatever limitations there are on
the predictability of the task, we are unable to find other classifiers that
have lower cost with respect to $\fairconstname$.
%
We allow $h_1$ and
$h_2$ to be learned in any way, as long as they are calibrated. Without loss of
generality, for the remainder of the paper, we will assume that
$\fairconst{h_1}{1} \ge \fairconst{h_2}{2}$.

Since by assumption we have no way to find a classifier for $G_1$ with lower
cost than $h_1$,
our goal is therefore to find
a classifier $\fair h_2$ with cost equal to $h_1$. This pair of classifiers
would represent the lowest cost (and therefore optimal) set of classifiers that
satisfies calibration and the equal cost constraint.
For a given base rate $\mu_t$ and value of the cost
function $\fairconstname$, a calibrated classifier's position in the
generalized false-positive/false-negative plane is uniquely determined
(\autoref{fig:levelorder}).
This is because each level-order
curve of the cost function $\fairconstname$ has negative
slope in this plane, and each level order curve only intersects a group's calibrated
classifier line once.
In other words, there is a unique solution in the false-positive/false-negative
plane for classifier $\fair h_2$ (\autoref{fig:postprocess}).

Consider the range of values that $\fairconstname$ can take. As noted above,
$\fairconst{h_t}{t} \ge 0$, with equality if and only if $h_t$ is the perfect
classifier. On the other hand,
the trivial classifier
(again, which outputs the constant $\br_t$ for all inputs)
is the calibrated classifier that achieves maximum cost for any $\fairconstname$
(see \autoref{lma:trivial} in \autoref{app:cost-funcs}).
As a result, the cost of a classifier for group
$G_t$ is between 0 and $\fairconst{h^{\br_t}}{t}$.
This naturally leads to a characterization of feasibility: \autoref{def:fair}
can be achieved if and only if $h_1$ incurs less cost than group $G_2$'s
trivial classifier $h^{\br_2}$; i.e.
if $\fairconst{h_1}{1} \leq \fairconst{h^{\mu_2}}{2}$.
This can be seen graphically in Figure \ref{fig:infeasible-trivial},
in which the level-order curve for $\fairconst{h_1}{1}$ does not intersect the set
of calibrated classifiers for $G_2$. Since, by assumption, we cannot find a
calibrated classifier for $G_1$ with
strictly smaller cost than $h_1$, there is no feasible solution.
On the other hand, if $h_1$ incurs less cost than $h^{\br_2}$, then we
will show feasibility by construction with a simple algorithm.

\paragraph{An Algorithm.}
While it may be possible to encode the constraints of \autoref{def:fair} into the training
procedure of $h_1$ and $h_2$, it is not immediately obvious how to do so.
Even naturally probabilistic algorithms, such as logistic regression,
can become uncalibrated in the presence of optimization constraints (as
is the case in \cite{zafar2017fairness}). It is not straightforward
to encode the calibration constraint if the probabilities
are assumed to be continuous, and post-processing calibration methods
\cite{platt1999probabilistic,zadrozny2001obtaining}
would break equal-cost constraints by modifying classifier scores.
Therefore, we look to achieve the calibrated Equalized Odds relaxation
by post-processing existing calibrated classifiers.

Again, given $h_1$ and $h_2$ with $\fairconst{h_1}{1} \ge \fairconst{h_2}{2}$, we want
to arrive at a calibrated classifier $\fair h_2$ for group $G_2$ such that
$\fairconst{h_1}{1} = \fairconst{\fair h_2}{2}$. Recall that, under our
assumptions, this would be
the best possible solution with respect to classifier cost. We show that this cost
constraint can be achieved by withholding predictive information for a randomly
chosen subset of group $G_2$. In other words, rather than always returning $h_2(\x)$
for all samples, we will occasionally return the group's mean probability
(i.e. the output of the trivial classifier $h^{\mu_2}$). In
\autoref{lma:interpolation} in \autoref{app:cost-funcs}, we show
that if
\begin{equation}
  \label{eqn:interpolation}
  \fair h_2(\x) = \begin{cases}
    h^{\br_2}(\x) = \br_2 & \text{with probability } \alpha \\
    h_2(\x) & \text{with probability } 1-\alpha
  \end{cases}
\end{equation}
then the cost of $\fair h_2$ is a linear interpolation between the costs of
$h_2$ and $h^{\br_2}$ (\autoref{fig:mixing}). More formally, we
have that $\fairconst{\fair h_2}{2} = (1-\alpha)
\fairconst{h_2}{2} + \alpha \fairconst{h^{\br_2}}{2}$), and thus setting $\alpha =
\frac{\fairconst{h_1}{1} - \fairconst{h_2}{2}}{\fairconst{h^{\br_2}}{2} -
\fairconst{h_2}{2}}$ ensures that $\fairconst{\fair h_2}{2} =
\fairconst{h_1}{1}$ as desired (\autoref{fig:postprocess}).
Moreover, this randomization preserves calibration (see
\autoref{sec:optimality-proofs}).
\autoref{alg:method} summarizes this method.
\begin{algorithm}
  \caption{Achieving Calibration and an Equal-Cost Constraint via Information Withholding}
  \label{alg:method}
  \begin{algorithmic}
    \STATE {\bfseries Input:} classifiers $h_1$ and $h_2$ s.t.
    $\fairconst{h_2}{2} \leq \fairconst{h_1}{1} \leq \fairconst{h^{\br_2}}{2}$, holdout set $P_{valid}$.

    \begin{itemize}
      \STATE Determine base rate $\br_2$ of $G_2$ (using
        $P_{valid}$) to produce trivial classifier $h^{\br_2}$.

      \STATE Construct $\fair h_2$ using with
        $ \alpha = \frac{\fairconst{h_1}{1} - \fairconst{h_2}{2}}{\fairconst{h^{\br_2}}{2} - \fairconst{h_2}{2}}$,
        where $\alpha$ is the interpolation parameter.
    \end{itemize}

    \RETURN $h_1$, $\fair h_2$ --- which are calibrated and satisfy $\fairconst{h_1}{1} = \fairconst{\fair h_2}{2}$.
  \end{algorithmic}
\end{algorithm}

%Note that \autoref{alg:method} is not the only possible post-processing method that achieves \autoref{def:fair}.
%However, any such post-processing method inherently relies on randomization. Any deterministic post-processing
%method (other than the identity function or returning the trivial classifier) would break the
%calibration constraint, as the post-processed scores would not match true probabilities.

\paragraph{Implications.}
In a certain sense, \autoref{alg:method} is an ``optimal'' method
because it arrives at the unique false-negative/false-positive solution for $\fair h_2$,
where $\fair h_2$ is calibrated and has cost equal to $h_1$. Therefore
(by our assumptions) we can find no better classifiers that satisfy \autoref{def:fair}.
%\gp{I cut the lower bound sentence}
%One might view this notion of ``optimality'' as a
%lower bound for the performance that must be sacrificed to satisfy
%\autoref{def:fair}. \gp{Maybe cut?}
%
%If, by assumption, we are
%unable to find a classifier for $G_1$ with lower cost than $h_1$ (and therefore
%lower false-positive or false-negative rates), then it must
%be the case that the final classifiers $h_1$ and $\fair h_2$ must have
%cost at least $\fairconst{h_1}{1}$. \autoref{alg:method} achieves this with
%equality, such that $\fairconst{\fair h_2}{2} = \fairconst{h_1}{1} =
%\fairconst{h_1}{1}$.
This simple result has strong consequences, as the tradeoffs to
satisfy both calibration and the equal-cost constraint are often unsatisfactory
--- both intuitively and experimentally (as we will show in \autoref{sec:experiments}).

We find two primary objections to this solution. First, it equalizes
costs simply by making a classifier strictly worse for one of the groups.
%It is not clear that this is desirable in many situations.
Second, it achieves this
cost increase by withholding information on a randomly chosen population subset,
making the outcome inequitable within the group (as measured by a standard measure
of inequality like the Gini coefficient). Due to the optimality of the
algorithm, the former of these issues is unavoidable in \emph{any} solution that
satisfies \autoref{def:fair}. The latter, however, is slightly more subtle, and
brings up the question of \emph{individual fairness} (what guarantees we would
like an algorithm to make with respect to each individual) and how it interacts
with \emph{group fairness} (population-level guarantees). While this certainly
is an important issue for future work, in this particular setting, even if one
could find another algorithm that distributes the burden of additional cost more
equitably, any algorithm will make at least as many false-positive/false-negative
errors as \autoref{alg:method}, and these misclassifications will always be
tragic to the individuals whom they affect.
The performance loss across the entire group is often
significant enough to make this combination of constraints somewhat worrying to
use in practice, regardless of the algorithm.

%Second, it suffers from
%a slightly more subtle flaw:
%the equal-cost guarantees it makes are in
%expectation, but the actual realized outcomes will be inequitable.
%To be more
%specific, consider a measure of inequality, like the Gini coefficient, which
%considers the expectation of a convex function to the outcome of each individual
%(as opposed to the expected outcome of a group, which is what the equal-cost
%constraint seeks to equalize). Intuitively, such a measure is small if outcomes
%are similar, but is large when they vary. The randomization of
%\autoref{alg:method} produces such varying outcomes: an $\alpha$ fraction of the
%population of $G_2$ gets a score of $\mu_2$, incurring the largest possible cost
%with respect to $\fairconstname$, while the remaining $1-\alpha$ fraction
%receives relatively little cost. The algorithm effectively achieves equal cost
%in expectation across groups by choosing a random subset of $G_2$ and saddling
%them with high cost, leading to inequity in outcomes.
%This inequality directly stems from the randomness of \autoref{alg:method},
%which, as stated earlier, is a property inherent to any calibrated post-processing approach.

%This particular objection regarding unequal outcomes is specific to the proposed
%algorithm; however, because it is optimal with repsect to performance, the
%objection that a certain amount of performance must be sacrificed to satisfy the
%equal-cost constraint holds for \emph{any} calibrated classifier that satisfies
%the equal-cost constraint. This means that this particular formulation of
%non-discrimination can only be achieved at the cost of explicitly reducing the
%classifier's accuracy on a particular group, regardless of what algorithm is
%used. As we will show in \autoref{sec:experiments}, we will show that the
%performance sacrifice incurred by this algorithm (and therefore any algorithm)
%can be rather large, even in expectation \gp{Does this belong?}.

%Assume we are given calibrated classifiers $h_1$ and $h_2$ for groups $G_1$ and
%$G_2$, trained through any procedure. We assume
%that these are the ``best'' classifiers we can find for each group with respect
%to the cost function $\fairconstname$ -- otherwise, we would start of with a
%better classifier instead.
%Without loss of generality, for the remainder of the paper we assume
%that \fairconst{h_1}{1} is larger than \fairconst{h_2}{2}. Figure
%\ref{newfigure} illustrates under what conditions \methodname{} can be
%satisfied. Intuitively, \methodname{} is achievable if and only if the given
%classifier $h_1$ has lower cost than the worst possible classifier for $G_2$. To
%see why this is, assume that this is not the case, that is, $h_1$ has higher
%cost than the worst possible classifier for $G_2$. Then, it is trivially
%impossible to derive new classifiers with the same costs for each group, simply
%because doing so would imply that we can find a classfier for $G_1$ with lower
%cost than $h_1$, which violates our original assumption.
%In the next sections, we will show that, under minor assumptions, \methodname{} is
%a {\bf feasible} notion of non-discrimination. Furthermore, we prove that we can
%derive {\bf optimal} non-discriminatory classifiers from (discriminatory)
%classifiers via a simple post-processing method. Finally, we show that
%\methodname{} is the {\bf strongest} achievable notion of non-discrimination
%that preserves calibration.
%\mr{Not sure if we still want this paragraph in light of the changes we're
%making to the remaining sections}
